\appendix
\titleformat{\chapter}[display]{\sffamily\bfseries\color{navy}}{\large APPENDIX \thechapter}{4pt}{\Huge}[\vspace{2pt}{\color{orange}\titlerule[2pt]}]
\renewcommand{\chaptermark}[1]{\markboth{Appendix \thechapter: #1}{}}

\chapter{Answers to Try It Yourself}\label{app:answers}
\small
\begin{multicols}{2}
\textbf{1.1} $343\times1/2=171.5$ m.\\
\textbf{1.2} IMU; time synchronization; lever arm and boresight; trajectory estimation (GNSS/INS + SLAM); ground control and checkpoints.

\textbf{2.1} $8.727\times10^{-4}$ rad; $80\times8.727\times10^{-4}=6.98$ cm.\\
\textbf{2.2} $\sqrt{4+2.25+9}=\sqrt{15.25}=3.91$.\\
\textbf{2.3} $\sqrt{144+25}=13$ m.\\
\textbf{2.4} 5 steps (a 3-4-5 triangle).

\textbf{3.1} $2.997\times10^8\times166.8\times10^{-9}/2=25.0$ m.\\
\textbf{3.2} $\delta t=2\times0.01/c\approx66.7$ ps.

\textbf{4.1} $0.008+40\times0.0005=0.028$ m $=2.8$ cm.\\
\textbf{4.2} $(60/20)^2=9$ times weaker.\\
\textbf{4.3} The white sign: it reflects much more light (higher reflectance); black rubber absorbs it.\\
\textbf{4.4} The tree (through gaps between leaves) and the glass storefront (the laser mostly passes through clear glass). Not the brick wall (opaque) or the pond (near-infrared light stops at the water surface).

\textbf{5.1} $\rho\cos V=39.848$; $x=34.509$, $y=-19.924$, $z=3.486$ m; $\sqrt{34.509^2+19.924^2+3.486^2}=40.000$ \checkmark\\
\textbf{5.2} $\rho=\sqrt{9+16+144}=13.0$ m; $H=\operatorname{atan2}(4,3)=53.13^\circ$; $V=\operatorname{atan2}(12,5)=67.38^\circ$.\\
\textbf{5.3} $32\times900\times20=576\,000$ pts/s; $10/20=0.5$ m per head line.

\textbf{6.1} $(0,10)$ then back to $(10,0)$: a rotation by $-\psi$ undoes $+\psi$ ($\vect R^{\mathsf T}=\vect R^{-1}$).\\
\textbf{6.2} $\psi=90^\circ-135^\circ=-45^\circ$.\\
\textbf{6.3} $0.02/40=5\times10^{-4}$ rad $=0.029^\circ$ (about 1.7 arc-minutes).

\textbf{7.1} $10\,\text{m}g=0.098$ m/s$^2$; $\tfrac12\times0.098\times10^2=4.9$ m.\\
\textbf{7.2} $10^\circ/\text{h}=4.85\times10^{-5}$ rad/s; $\tfrac12\times20\times4.85\times10^{-5}\times30^2=0.44$ m.\\
\textbf{7.3} $9.83-9.80=0.03$ m/s$^2\approx3$ m$g$.

\textbf{8.1} $30\times0.0005=1.5$ cm.\\
\textbf{8.2} $90^\circ/\text{s}\times0.002=0.18^\circ=0.00314$ rad; $\times15=4.7$ cm.\\
\textbf{8.3} The car moves while the picture is being taken, so different parts are captured at different positions --- just like a LiDAR sweep.

\textbf{9.1} 0.765 m error in the body $z$ direction; it does \emph{not} depend on range (a lever-arm error is a constant offset in the body frame).\\
\textbf{9.2} $0.02/60=3.3\times10^{-4}$ rad $=0.019^\circ$.\\
\textbf{9.3} $Z=0.974$ m, so the point height is 10.974 m (31 cm lower than with roll and pitch).

\textbf{10.1} Straight line: $(10.40,\ 5.00,\ 0^\circ)$.\\
\textbf{10.2} $K=0.09/0.18=0.5$; $\mu=50.10$ m; $\sigma^2=0.045$, $\sigma=0.212$ m.\\
\textbf{10.3} The landmark is somewhere known, so seeing it at a measured distance and direction rules out many possible positions.

\textbf{11.1} Centered: $\vect p'=(\mp0.5,0)$, $\vect q'=(0,\mp0.5)$; sums $0$ and $0.5$ $\Rightarrow\theta=90^\circ$; $\vect t=(2,1.5)-\vect R(0.5,0)=(2,1.0)$.\\
\textbf{11.2} The line is the $z$-axis: $d=\sqrt{3^2+4^2}=5$.\\
\textbf{11.3} Everything looks the same as you move, so there is nothing to tell you how far you have gone.

\textbf{12.1} Odometry sums to exactly 3.0, so there is no misclosure: $x=(1.0,\,2.1,\,3.0)$, all residuals 0.\\
\textbf{12.2} $\begin{psmallmatrix}2&-1&0\\-1&2&-1\\0&-1&5\end{psmallmatrix}\vect x=\begin{psmallmatrix}-0.1\\0.2\\13.7\end{psmallmatrix}$ $\Rightarrow x=(1.062,\ 2.223,\ 3.185)$; odometry residuals $+0.062$ each, GPS $-0.015$. The stiffer GPS gets most of its way.\\
\textbf{12.3} The stiff band hardly stretches; the others stretch more to share the strain.

\textbf{13.1} Mean $1.2$ cm; RMSE $\sqrt{34/5}=2.61$ cm; with checkpoint error $\sqrt{6.8+2.25}=3.01$ cm.\\
\textbf{13.2} $\sqrt{1+4+4+16}=5$ cm; attack the 4 cm term first (it is 64\% of the variance).\\
\textbf{13.3} If the computer used them to fit the map, they would of course agree with it --- the test would prove nothing.

\textbf{14.1} Along the beam: 0.3 cm; across: $15\times1.745\times10^{-4}=0.26$ cm.\\
\textbf{14.2} Fully correlated errors add linearly: $2+3=5$ cm (not 3.6 cm).\\
\textbf{14.3} $0.5/\sqrt n<0.05\Rightarrow n>100$.\\
\textbf{14.4} $1.96\times1.5=2.94$ cm.

\textbf{15.1} $8\times0.01227=0.098$ m (9.8 cm).\\
\textbf{15.2} $0.005+3\times0.00611=0.023$ m (2.3 cm).\\
\textbf{15.3} $\sqrt{1.25^2+0.65^2}=1.41$ cm.\\
\textbf{15.4} $1800\times20\times10^{-6}=0.036$ s; $\times10$ m/s $=0.36$ m.

\textbf{16.1} Misclosure 0 (a constant bias rotates the whole loop); far corner $707\times0.008727=6.2$ m off.\\
\textbf{16.2} $0.01\times60=0.6^\circ/\sqrt{\text{h}}$.\\
\textbf{16.3} $0.03/25=0.0012$ rad $=0.069^\circ$.\\
\textbf{16.4} Because an error that is the same on every leg rotates or stretches the whole figure, and a rotated or stretched closed figure still closes.\\
\textbf{16.5} $b=0.0065/1.745\times10^{-3}=3.7$ m.\\
\textbf{16.6} $\sqrt{0.30^2+0.022^2+0.321^2}=0.44$ m.

\textbf{17.1} $152.4\times\sin37.5^\circ\times1.745\times10^{-4}=0.016$ m (1.6 cm).\\
\textbf{17.2} Zone depth $c/(2\times1.2\times10^6)=124.9$ m; $152.4/124.9=1.22$ $\Rightarrow$ zone 2.\\
\textbf{17.3} $\sqrt{1.0^2+2.5^2+0.5^2}=2.7$ cm.\\
\textbf{17.4} It is Class 3R: from a moving aircraft the beam sweeps past any eye too fast to harm it (NOHD 0 m above 1 kn), but a stationary beam has a hazard distance (0.3 m, 3.5 m with optical aids).

\textbf{18.1} $(\lfloor40.52\rfloor,\lfloor-9.6\rfloor,\lfloor7.96\rfloor)=(40,-10,7)$.\\
\textbf{18.2} $16\times1\times0.75=12$ m$^3$.\\
\textbf{18.3} A roll boresight error (it tilts the road in opposite senses on opposite passes).

\textbf{19.1} 5J-17.051(2): the accuracy of data shall be independently verified. A LiDAR survey satisfies it by testing against independent checkpoints that were never used as control, and documenting the method and RMSE.\\
\textbf{19.2} March 5 --- the survey date is the last date of data acquisition (5J-17.051(3)(d)), not the processing date.\\
\textbf{19.3} No. The edited tile's hash no longer matches the code listed in the sealed signature file, so that file is no longer sealed (5J-17.062(4)).\\
\textbf{19.4} U.S.\ survey foot: 365\,760.7315 m; international foot: 365\,760.0000 m; difference 0.7315 m (2.40 ft).

\textbf{20.1} Exercise responsible charge: obtain and review the calibration, trajectory quality and control; run (or direct) an independent checkpoint test; confirm datums, units and classification; state responsibility for the vendor's data before sealing.\\
\textbf{20.2} Any three of: no data behind solid objects or under dense vegetation; no near-infrared returns below water; no underground utilities; occlusions by traffic; reduced accuracy beyond the usable range; areas affected by GNSS outages.\\
\textbf{20.3} Part 107 limits altitude to 400 ft AGL (absent a waiver or nearby structure); lower altitude also shortens the range, so attitude and range terms shrink and point density rises.
\end{multicols}
\normalsize

\chapter{Final Exam (Administered Separately)}
The final exam and its answer keys are \textbf{not printed in this book}. They are in the separate \emph{Instructor Packet}, which contains:
\begin{itemize}
\item \textbf{two exam versions} (A and B) of 50 conceptual and standards questions each --- 15 true/false, 20 multiple choice on concepts, and 15 on standards and professional practice;
\item the answer keys, a map of every question to the seminar's learning objectives, and administration instructions;
\item the \textbf{pass mark: 70\%} (35 of 50 correct).
\end{itemize}
No calculations are required and no calculator is needed. To prepare, review the figures, Simple Approach boxes, Big Ideas, Field Tips, and Modules 13, 14, 19 and 20.

\chapter{Equation Summary Sheet}
\small
\renewcommand{\arraystretch}{1.6}
\renewcommand{\arraystretch}{1.45}
\begin{longtable}{>{\raggedright\arraybackslash}p{4.1cm}>{\raggedright\arraybackslash}p{10.2cm}>{\centering\arraybackslash}p{1.1cm}}\toprule
\textbf{Quantity} & \textbf{Equation} & \textbf{Mod.}\\\midrule
\endfirsthead
\textbf{Quantity} & \textbf{Equation} & \textbf{Mod.}\\\midrule
\endhead

Time-of-flight range & $\rho=c\,\Delta t/2,\ \ c=c_0/n$ & 3\\
Max.\ unambiguous range & $\rho_{\max}=c/(2f_{\text{PRF}})$ & 3\\
Phase range & $\rho=\tfrac{\lambda_m}{2}(N+\Delta\varphi/2\pi)$ & 3\\
Beam footprint & $D=d_0+\rho\gamma$ & 4\\
Received power & $P_r\propto\rho_{\text{refl}}\cos\theta_i/R^2$ & 4\\
Polar to Cartesian & $x=\rho\cos V\cos H,\ y=\rho\cos V\sin H,\ z=\rho\sin V$ & 5\\
Point spacing & $\Delta_{\text{along}}=v/f,\ \Delta_{\text{across}}=\rho\Delta\theta$ & 5\\
Small-angle error & $\delta=\rho\,\delta\alpha$ (rad) & 2, 6\\
3-D rotation & $\vect R=\vect R_z(\psi)\vect R_y(\theta)\vect R_x(\phi)$ & 6\\
Rigid transform & $\vect r_W=\vect R\vect r_B+\vect p$;\ $\Delta\vect T_{ij}=\vect T_i^{-1}\vect T_j$ & 6\\
IMU model & $\hat{\boldsymbol\omega}=\boldsymbol\omega+\vect b^\omega+\vect n^\omega,\ \hat{\vect a}=\vect R^{BW}(\vect a-\vect g)+\vect b^a+\vect n^a$ & 7\\
Strapdown & $\vect v'=\vect v+\vect g\Delta t+\vect R\hat{\vect a}\Delta t$;\ $\vect p'=\vect p+\vect v\Delta t+\tfrac12(\vect g+\vect R\hat{\vect a})\Delta t^2$ & 7\\
Bias growth & $\delta y=\tfrac12vb_gt^2$;\ $\delta x=\tfrac12b_at^2$ & 7\\
Preintegration & $\Delta\vect R_{ij}=\vect R_i^{\mathsf T}\vect R_j$, \ldots & 7\\
Timing error & $\delta=v\,\delta t+\rho\dot\psi\,\delta t$ & 8\\
Direct georeferencing & $\vect r_W=\vect p_W+\vect R^W_B(\vect R^B_S\vect r_S+\vect l^B)$ & 9\\
Motion model & $\alpha=(R-L)/w$;\ $x'=x+(r+\tfrac w2)(\sin(\theta+\alpha)-\sin\theta)$ & 10\\
Kalman (1-D) & $K=\bar\sigma^2/(\bar\sigma^2+\sigma_z^2)$;\ $\mu=\bar\mu+K(z-\bar\mu)$;\ $\sigma^2=(1-K)\bar\sigma^2$ & 10\\
EKF & $\bar{\vect P}=\vect G\vect P\vect G^{\mathsf T}+\vect Q$;\ $\vect K=\bar{\vect P}\vect H^{\mathsf T}(\vect H\bar{\vect P}\vect H^{\mathsf T}+\vect R)^{-1}$ & 10\\
ICP & $\min\sum\|\vect R\vect p_i+\vect t-\vect q_i\|^2$;\ $\vect R=\vect V\vect U^{\mathsf T}$, $\vect t=\bar{\vect q}-\vect R\bar{\vect p}$ & 11\\
Edge / plane distance & $d_e=\|(\vect p-\vect a)\times(\vect p-\vect b)\|/\|\vect a-\vect b\|$; $d_p=|(\vect p-\vect u)\cdot\vect n|/\|\vect n\|$ & 11\\
Factor graph & $\vect X^*=\arg\min\sum\|\vect h_k(\vect X)-\vect z_k\|^2_{\Sigma_k}$ & 12\\
Normal equations & $\vect A^{\mathsf T}\vect W\vect A\,\vect x=\vect A^{\mathsf T}\vect W\vect b$ & 12\\
RMSE & $\sqrt{\tfrac1n\sum e_i^2}$;\ $\text{RMSE}_H=\sqrt{\text{RMSE}_x^2+\text{RMSE}_y^2}$ & 2, 13\\
Final accuracy & $\sqrt{\text{RMSE}_{\text{data}}^2+\text{RMSE}_{\text{survey}}^2}$ & 13\\
Error budget (RSS) & $\sigma=\sqrt{\sum\sigma_j^2}$ & 2, 13\\
Outlier test & $\bar d_i>\mu+\alpha\sigma$ & 18\\
Grid volume & $V=\sum A(z_{\text{top}}-z_{\text{base}})$ & 18\\
Variance propagation & $\sigma_y^2=\sum(\partial f/\partial x_i)^2\sigma_i^2+2\sum a_ia_j\sigma_{ij}$ & 14\\
Covariance propagation & $\boldsymbol\Sigma_y=\vect J\boldsymbol\Sigma_x\vect J^{\mathsf T}$ & 14\\
Averaging with bias & $\text{RMSE}=\sqrt{b^2+\sigma^2/n}$ & 14\\
95\% factors & 1.960 (1-D), 2.448 (2-D), 2.796 (3-D) & 14\\
Quantization & $\sigma_q=q/\sqrt{12}$ & 15\\
Angular step & $\Delta\theta=360^\circ/N_{\text{col}}$;\ $s=\rho\Delta\theta$ & 15\\
Random walks & $\text{ARW}=60N_g$;\ $\sigma_\psi=\text{ARW}\sqrt t$ & 16\\
Dual-antenna heading & $\sigma_\psi\approx\sigma_{\text{rel}}/b$ & 16\\
Coast error & $\sqrt{(\sigma_vt)^2+(\tfrac12b_at^2)^2+(\tfrac16gb_gt^3)^2}$ & 16\\\bottomrule
\end{longtable}
\normalsize

\chapter{Hands-On Labs}
These five labs turn the modules into things you can do with your hands, a tape measure, a spreadsheet or a laptop. Labs 1--2 use the simple approach and need no math; Labs 3--5 are for engineers.

\section*{Lab 1 --- Be a human LiDAR (Modules 3 and 5)}
\begin{kid}
\textbf{You need:} a protractor taped flat on a chair, a tape measure, graph paper, a friend.\\
\textbf{Do it:} Sit at the chair --- you are the scanner. Point at a corner of the room, a door edge, a table leg. For each, read the angle on the protractor (that is $H$) and have your friend measure the distance (that is $\rho$). Write them in a table.\\
\textbf{Then:} On graph paper, put yourself in the middle. For each object, turn your ruler to the angle and mark a dot at the distance (1 square = 50 cm). Join the dots --- you have drawn a \emph{point-cloud map} of your room!\\
\textbf{Engineer's extension:} compute each dot with $x=\rho\cos H$, $y=\rho\sin H$ and compare with the drawing.
\end{kid}

\section*{Lab 2 --- The drifting blindfold walk (Modules 7 and 10)}
\begin{kid}
\textbf{Do it:} Put a cone at the start. Blindfolded, try to walk 10 steps forward, turn around, and walk 10 steps back to the cone. Mark where you stop. Repeat three times. You will miss the cone --- that is \emph{drift}. Now try again, but you may touch one wall halfway (a \emph{landmark}). Do you get closer?\\
\textbf{What you learned:} counting steps is like an IMU: little mistakes add up. Touching something you know is like a LiDAR match or a loop closure: it fixes the mistakes.
\end{kid}

\section*{Lab 3 --- ICP in 30 lines of Python (Module 11)}
Run this with Python 3 and NumPy. It builds an L-shaped ``room,'' makes a sparse, moved and turned copy (a new scan), and recovers the motion with ICP. Watch the RMS gap fall each iteration --- exactly Figure 11.1.
\begin{tcolorbox}[base,colback=gray!4,colframe=grayx,title={icp\_lab.py}]
\scriptsize
\begin{verbatim}
import numpy as np

def best_fit(P, Q):
    """Closed-form rigid fit (Module 11): returns R, t with Q ~ R P + t."""
    pb, qb = P.mean(0), Q.mean(0)
    W = (P - pb).T @ (Q - qb)
    U, S, Vt = np.linalg.svd(W)
    R = Vt.T @ U.T
    if np.linalg.det(R) < 0:          # avoid a mirror image
        Vt[-1] *= -1; R = Vt.T @ U.T
    return R, qb - R @ pb

def icp(P, Q, iters=30):
    cur = P.copy()
    for k in range(iters):
        d = ((cur[:, None, :] - Q[None, :, :])**2).sum(-1)
        match = Q[d.argmin(1)]                 # 1. nearest neighbors
        R, t = best_fit(cur, match)            # 2. best rigid fit
        cur = cur @ R.T + t                    # 3. apply
        rms = np.sqrt(d.min(1).mean())
        print(f"iter {k:2d}  RMS gap = {rms:.4f} m")
        if rms < 1e-4: break                   # 4. repeat until stable
    return best_fit(P, cur)                    # total transform

def room(n):                        # an L-shaped wall plus a short wall
    s = np.linspace(0, 1, n)
    return np.vstack([np.c_[6*s, 0*s], np.c_[6+0*s, 3*s], np.c_[3.5+0*s, 3+1.5*s]])
Q = room(600)                       # dense "map"
th = np.radians(6); Rt = np.array([[np.cos(th), -np.sin(th)], [np.sin(th), np.cos(th)]])
P = (room(40) - [0.25, -0.15]) @ Rt   # sparse new scan, moved and turned
R, t = icp(P, Q)
print("recovered angle (deg):", np.degrees(np.arctan2(R[1, 0], R[0, 0])).round(3))
\end{verbatim}
\end{tcolorbox}
\textbf{Expected output:} RMS gap 0.18 m at iteration 0 falling to about 0.002 m; recovered angle $6.0^\circ$. \textbf{Experiments:} (a) increase the turn to 25° and watch ICP get stuck in a local minimum --- the reason it needs an IMU guess; (b) delete the short wall from \texttt{room()} and add a shift along $x$ only: the long walls let the scan slide --- the hallway effect.

\section*{Lab 4 --- Field check of boresight from two passes (Module 9)}
\begin{eng}
\textbf{Procedure:} pick a sharp vertical feature beside the road --- a building corner or a pole. Drive past it in both directions in the same lane. In the cloud, measure the along-track offset $\Delta$ between the feature's position in the two passes, and the range $\rho$ from the drive line to the feature.

\textbf{Why it works:} a yaw boresight error $\delta\psi$ shifts the feature by $\rho\,\delta\psi$ forward on one pass and backward on the other, so the passes differ by $2\rho\,\delta\psi$:
\[ \delta\psi\approx\frac{\Delta}{2\rho}. \]
\textbf{Example:} $\Delta=4.2$ cm along-track at $\rho=30$ m $\Rightarrow\delta\psi=0.042/60=7.0\times10^{-4}$ rad $=0.040^\circ$. The same logic applied to the \emph{height} of a flat road edge on both sides gives the roll residual, and to the height of features ahead/behind gives pitch.
\end{eng}

\section*{Lab 5 --- Checkpoint report in a spreadsheet (Module 13)}
\begin{eng}
Set up columns: \textsf{Point ID | Survey Z | LiDAR Z | dZ = LiDAR $-$ Survey | dZ$^2$}. Then:
\begin{itemize}
\item mean: \texttt{=AVERAGE(D2:D31)}
\item standard deviation: \texttt{=STDEV.S(D2:D31)}
\item RMSE: \texttt{=SQRT(SUMSQ(D2:D31)/COUNT(D2:D31))}
\item final RMSE with checkpoint survey error in G1: \texttt{=SQRT(RMSE\^{}2+G1\^{}2)}
\end{itemize}
Enter the ten residuals from the Module 13 worked example and confirm 0.53 cm, 2.64 cm, 2.56 cm and 2.75 cm. Then add a deliberate blunder (+25 cm) to one point and see how strongly RMSE reacts --- and why each large residual must be investigated, not just averaged.
\end{eng}

\chapter{Sample SOP and Client Accuracy Agreement}
This appendix is a \textbf{template}. It shows how the ideas in this course become a written Standard Operating Procedure (SOP) and a signed agreement on horizontal and vertical accuracy. Replace every sample value with your project's requirements, and have the surveyor in responsible charge review it against the current standards and contract.

\begin{center}\small
\renewcommand{\arraystretch}{1.25}
\begin{tabularx}{\linewidth}{>{\bfseries}p{3.4cm}X>{\bfseries}p{2.6cm}X}\toprule
SOP number & ML-SOP-001 & Revision & 0 (sample)\\
Title & \multicolumn{3}{p{10.5cm}}{Mobile LiDAR Topographic Survey --- Collection, Processing, QA/QC and Reporting}\\
Effective date & \_\_\_\_\_\_\_\_ & Approved by & \_\_\_\_\_\_\_\_, PSM\\\bottomrule
\end{tabularx}\end{center}

\section*{1. Purpose}
To produce mobile LiDAR point clouds and derived topographic products whose horizontal and vertical accuracy is agreed with the client \emph{before} collection, measured independently \emph{after} processing, and reported in the language of the ASPRS Positional Accuracy Standards.

\section*{2. Scope}
Vehicle-, backpack- or cart-mounted mobile laser scanning for topographic surveys, corridor mapping, as-builts and volumes. It does not cover boundary determination or subsurface utility locating.

\section*{3. References}
\begin{itemize}
\item ASPRS Positional Accuracy Standards for Digital Geospatial Data, Edition 2, Version 2 (2024).
\item ASPRS LAS Specification, version 1.4 (current revision).
\item Chapter 5J-17, Florida Administrative Code, Standards of Practice (current text).
\item FDOT guidance for terrestrial mobile LiDAR and the FDOT Florida Permanent Reference Network (FPRN), where applicable.
\item IEC 60825-1, Safety of laser products (Class 1 equipment only).
\item Manufacturer manuals and datasheets for the scanner, IMU/INS and GNSS receivers.
\end{itemize}

\section*{4. Roles and responsibilities}
\begin{center}\small
\renewcommand{\arraystretch}{1.2}
\begin{tabularx}{\linewidth}{>{\bfseries}p{3.6cm}X}\toprule
Surveyor in responsible charge (PSM) & approves the SOP and accuracy agreement; signs the final report and accuracy statement\\
Project manager & client contact; schedule; scope changes; nonconformance decisions with the PSM\\
Field crew chief & equipment checks, calibration loop, collection, field log, safety\\
Control surveyor & control and checkpoint survey; keeps checkpoint coordinates \emph{sealed} from processing staff\\
Processing analyst & trajectory, georeferencing, cleaning, classification, products\\
QA reviewer (independent) & accuracy testing against checkpoints, swath checks, report review\\\bottomrule
\end{tabularx}\end{center}

\section*{5. Client accuracy agreement}
Complete and sign this agreement \textbf{before mobilization}. It fixes what will be delivered, how accurate it must be, and how accuracy will be proven.

\begin{center}\small
\renewcommand{\arraystretch}{1.3}
\begin{tabularx}{\linewidth}{>{\raggedright\arraybackslash}p{5.0cm}X}\toprule
\textbf{Item} & \textbf{Agreed value (sample --- edit for your project)}\\\midrule
Project / location & Main Street corridor, 2.0 km, curb to right-of-way line\\
Horizontal datum and realization & NAD 83(2011), epoch 2010.0; State Plane Coordinate System, Florida East zone\\
Vertical datum and geoid & NAVD 88 via GEOID18 (state the model)\\
Linear units & \_\_\_\_\_\_ (state U.S. survey foot or international foot explicitly; NGS deprecated the U.S. survey foot after 2022 --- confirm statutory and client requirements)\\
\textbf{Horizontal accuracy class} & \textbf{RMSE$_H$ $\le$ 5 cm} (0.16 ft) on well-defined features within 25 m (82 ft) of the trajectory\\
\textbf{Vertical accuracy class (NVA)} & \textbf{RMSE$_V$ $\le$ 3 cm} (0.10 ft) on hard, non-vegetated surfaces --- \emph{acceptance criterion}\\
Vegetated vertical accuracy (VVA) & tested and reported as found (no pass/fail)\\
Data internal precision & within-swath smooth-surface max difference $\le$ 1.8 cm (0.06 ft); swath-to-swath RMSD$_z$ $\le$ 2.4 cm (0.08 ft), max difference $\le$ 4.8 cm (0.16 ft) (0.6$X$ / 0.8$X$ / 1.6$X$ for $X$ = 3 cm)\\
Checkpoints & $\ge$30 NVA checkpoints on hard surfaces; $\ge$30 well-defined horizontal checkpoints (painted chevrons, corners); distributed along the full length and both sides; never used as control\\
Checkpoint / control accuracy & checkpoints RMSE$_V$ $\le$ 1.5 cm (0.05 ft), RMSE$_H$ $\le$ 2.5 cm (0.08 ft) (at least 2$\times$ better than the product); method and redundancy documented\\
Accuracy reporting & ASPRS Edition 2 wording; checkpoint survey error included in the reported RMSE\\
Point density & $\ge$ 400 points/m$^2$ (37 points/ft$^2$) on pavement within 20 m (66 ft) (verified on test patches)\\
Usable range & points beyond 25 m (82 ft) from the trajectory excluded from accuracy-critical products\\
Classification (LAS codes) & 1, 2, 6, 7, 11, 13--15, 17, 18 (others by agreement)\\
Deliverables & LAS 1.4 (PDRF 7, RGB) as LAZ tiles; ground surface (TIN) with breaklines; planimetric line work in CAD format; trajectory file; control and checkpoint report; survey report with accuracy statement; metadata\\
Collection conditions & dry pavement; daylight or as agreed; traffic control per permit; GNSS/RTK availability checked (e.g.\ FPRN VRS)\\
Acceptance and remedy & if NVA or horizontal RMSE exceeds the class, producer investigates, reprocesses, or recollects the affected area at no added cost; outliers beyond 3$\times$RMSE are investigated and documented, never silently removed\\\midrule
Client representative & Name \_\_\_\_\_\_\_\_\_\_\_\_ \quad Signature \_\_\_\_\_\_\_\_\_\_\_\_ \quad Date \_\_\_\_\_\_\\
Surveyor in responsible charge & Name \_\_\_\_\_\_\_\_\_\_\_\_ \quad PSM No.\ \_\_\_\_\_\_ \quad Signature \_\_\_\_\_\_\_\_\_\_\_\_ \quad Date \_\_\_\_\_\_\\\bottomrule
\end{tabularx}\end{center}

\begin{tip}
\textbf{Agree on what ``accuracy'' means before anyone drives.} Most disputes come from four gaps: (1) the class was stated without saying RMSE or 95\%; (2) horizontal and vertical were not separated; (3) nobody said which surfaces are tested (hard surfaces vs.\ grass); (4) datum, geoid and units were assumed. The table above closes all four.
\end{tip}

\section*{6. Procedure}
\subsection*{6.1 Planning}
\begin{enumerate}[label=\alph*.]
\item Build an error budget for the system and route (Modules 13--14); confirm the agreed classes are achievable at the planned range and speed.
\item Plan routes for loop closures and passes in both directions; identify GNSS-poor areas (tunnels, canopy, urban canyons) and add control there.
\item Lay out control (every $\sim$250 m / 820 ft and in GNSS-poor areas) and checkpoints (independent, distributed, both sides).
\item Confirm RTK source (base or network such as FPRN), mountpoint, datum and epoch.
\item Obtain permits and maintenance-of-traffic plans; schedule for dry surfaces.
\end{enumerate}

\subsection*{6.2 Equipment checks and calibration}
\begin{enumerate}[label=\alph*.]
\item Verify firmware, calibration file (lever arm, boresight) and its date; record in the field log.
\item Confirm time synchronization (PPS/PTP lock) and IMU, LiDAR and GNSS streams at nominal rates.
\item Static IMU alignment and a calibration loop over known features; review boresight residuals before production.
\item Laser safety: street and terrestrial scanners Class 1; airborne Class 3R scanners (e.g.\ RIEGL VUX-240\textsuperscript{24}) operated only per the manufacturer's NOHD/ENOHD and the company laser-safety program; post crew safety briefing.
\end{enumerate}

\subsection*{6.3 Control and checkpoint survey}
\begin{enumerate}[label=\alph*.]
\item Survey control and checkpoints with methods meeting the agreed accuracy (redundant RTK, static GNSS, or total station / digital levels).
\item Assign IDs; deliver checkpoint coordinates only to the QA reviewer (sealed from processing).
\item Check into a known monument at start and end of each session; record differences.
\end{enumerate}

\subsection*{6.4 Data collection}
\begin{enumerate}[label=\alph*.]
\item Collect at or below the planned speed; drive both directions; avoid stops in traffic where possible.
\item Keep a field log: times, weather, GNSS outages, sensor warnings, deviations from plan.
\item Verify data on site (coverage, density, no gaps) before demobilizing.
\end{enumerate}

\subsection*{6.5 Processing}
\begin{enumerate}[label=\alph*.]
\item Trajectory: GNSS/INS (RTK or PPK) with LiDAR SLAM and loop closures; review trajectory standard deviations and flag weak stretches.
\item Georeference with the approved calibration; de-skew; apply control.
\item Clean noise and moving objects; classify to the agreed LAS codes; document voxel size and filter settings.
\item Extract products: surface with breaklines, planimetrics, volumes as ordered.
\end{enumerate}

\subsection*{6.6 QA/QC and acceptance testing}
\begin{center}\small
\renewcommand{\arraystretch}{1.2}
\begin{tabularx}{\linewidth}{>{\raggedright\arraybackslash}p{4.2cm}Xc}\toprule
\textbf{Test} & \textbf{Criterion (sample)} & \textbf{Pass?}\\\midrule
NVA, $\ge$30 hard-surface checkpoints & RMSE$_V$ (incl.\ checkpoint error) $\le$ 3 cm (0.10 ft) & $\square$\\
Horizontal, $\ge$30 well-defined points & RMSE$_H$ (incl.\ checkpoint error) $\le$ 5 cm (0.16 ft) & $\square$\\
VVA & reported as found & $\square$\\
Within-swath precision & max difference $\le$ 1.8 cm (0.06 ft) on flat surfaces & $\square$\\
Swath-to-swath & RMSD$_z$ $\le$ 2.4 cm (0.08 ft); max $\le$ 4.8 cm (0.16 ft) & $\square$\\
Point density & $\ge$ 400 pts/m$^2$ on test patches & $\square$\\
Classification & spot check; no ground points on vehicles/buildings & $\square$\\
Datum / geoid / units & match the agreement and file headers & $\square$\\
Residual review & every residual $>$3$\times$RMSE investigated and documented & $\square$\\\bottomrule
\end{tabularx}\end{center}

\subsection*{6.7 Reporting}
The report states methods, equipment, calibration, trajectory solution, control used, checkpoints used (kept separate), datums, geoid, units, the test results and the accuracy statement. Sample statement:
\begin{tcolorbox}[base,colback=gray!4,colframe=grayx,title={Sample accuracy statement}]
``This data set was tested to meet ASPRS Positional Accuracy Standards for Digital Geospatial Data, Edition 2, Version 2 (2024) for a 3 cm RMSE$_V$ Vertical Accuracy Class and a 5 cm RMSE$_H$ Horizontal Accuracy Class. Non-Vegetated Vertical Accuracy (NVA) was found to be RMSE$_V$ = 2.1 cm using 30 independent checkpoints, including checkpoint survey uncertainty. Horizontal accuracy was found to be RMSE$_H$ = 2.8 cm using 30 well-defined checkpoints. Vegetated Vertical Accuracy (VVA) was found to be RMSE$_V$ = \_\_ cm. Horizontal datum: NAD 83(2011) epoch 2010.0, SPCS Florida East; vertical datum: NAVD 88 (GEOID18); units: \_\_\_.''
\end{tcolorbox}

\subsection*{6.8 Nonconformance and corrective action}
If any acceptance test fails: stop delivery of the affected area, identify the cause (trajectory, calibration, control, datum, processing), correct and retest, and record the event, cause and fix in the project file. Never delete checkpoints to make a test pass.

\subsection*{6.9 Records and retention}
Archive raw sensor data, trajectory, calibration file, control and checkpoint data, processing logs and settings (including voxel and filter parameters), QA results, deliverables and the signed agreement, for the retention period required by law and contract.

\section*{7. Revision history}
\begin{center}\small
\begin{tabular}{llll}\toprule
Rev. & Date & Description & Approved\\\midrule
0 & \_\_\_\_\_\_ & Initial issue (sample from course text) & \_\_\_\_\_\_\\\bottomrule
\end{tabular}\end{center}


\chapter{Glossary}
\small
\begin{multicols}{2}
\begin{description}[leftmargin=0pt,itemsep=2pt,font=\bfseries\color{navy}]
\item[Attitude] Orientation of the platform: roll, pitch, yaw.
\item[Bias] A slowly changing offset in a sensor reading; integrates into drift.
\item[Boresight] The small rotation between the LiDAR and IMU frames.
\item[Checkpoint] A surveyed point used only to test accuracy, never to adjust.
\item[Control point] A surveyed point used to fix the datum and constrain the solution.
\item[Covariance] A matrix describing uncertainties and how they are correlated.
\item[De-skew] Correcting each point of a sweep for the motion during that sweep.
\item[Dead reckoning] Computing position by adding up measured motion from a known start.
\item[DTM] Digital terrain model: the bare-earth surface.
\item[EKF] Extended Kalman filter: a Kalman filter linearized with Jacobians.
\item[Factor graph] A graph of unknown states and the measurements (factors) that connect them, solved by least squares.
\item[Footprint] Size of the laser spot on the target.
\item[GNSS] Global Navigation Satellite Systems (GPS, GLONASS, Galileo, BeiDou).
\item[Hallway effect] Weak along-track constraint in featureless corridors.
\item[ICP] Iterative Closest Point: repeatedly pair and fit two point clouds.
\item[IMU] Inertial measurement unit: accelerometers + gyroscopes.
\item[Innovation] Measurement minus its prediction.
\item[iSAM2] Incremental smoothing and mapping: updates only the changed part of a factor graph.
\item[Jacobian] Matrix of partial derivatives used to linearize a model.
\item[Kalman gain] Weight given to a new measurement versus the prediction.
\item[Keyframe] A selected scan whose pose becomes a node in the graph.
\item[LAS/LAZ] Standard point-cloud file format / its compressed version.
\item[Lever arm] Offset vector between two sensor origins.
\item[LIO-SAM] Tightly-coupled lidar--inertial odometry via smoothing and mapping (Shan \emph{et al.}, 2020).
\item[Loop closure] Recognizing a revisited place and adding that constraint.
\item[MAP] Maximum a posteriori: the most probable solution given all data.
\item[NDT] Normal Distributions Transform: registration against per-voxel Gaussians.
\item[Point cloud] A set of 3-D points, often with intensity and time.
\item[Pose] Position plus orientation (6 degrees of freedom).
\item[PPS] Pulse per second: GNSS timing signal for synchronization.
\item[Preintegration] Summarizing many IMU samples into one relative-motion measurement.
\item[RMSE] Root-mean-square error: the standard accuracy measure.
\item[RSS] Root-sum-square: how independent errors combine.
\item[SLAM] Simultaneous localization and mapping.
\item[Smoother] An estimator that uses past \emph{and} future data to refine each state.
\item[Strapdown] IMU fixed to the body; orientation tracked computationally.
\item[Tightly coupled] Sensors are fused in one estimation problem, at raw-measurement level.
\item[Voxel] A small 3-D grid cube.
\end{description}
\end{multicols}
\normalsize

\chapter{References and Further Study}
\small
\begin{enumerate}[leftmargin=*]
\item Shan, T., Englot, B., Meyers, D., Wang, W., Ratti, C., \& Rus, D. (2020). LIO-SAM: Tightly-coupled lidar inertial odometry via smoothing and mapping. \emph{IEEE/RSJ IROS}. arXiv:2007.00258. (Source document for Modules 7, 11, 12.)
\item \emph{Mobile LiDAR Basics for Surveying and Mapping} --- technical seminar document, October 2026 (59 pp.). (Source document; examples revised and corrected here.)
\item \emph{LiDAR Mobile Mapping with SLAM: Course Outline}, 2026. (Source document; defines the four-module structure mapped in the course map.)
\item Zhang, J., \& Singh, S. (2017). Low-drift and real-time lidar odometry and mapping (LOAM). \emph{Autonomous Robots}, 41(2), 401--416.
\item Forster, C., Carlone, L., Dellaert, F., \& Scaramuzza, D. (2017). On-manifold preintegration for real-time visual--inertial odometry. \emph{IEEE Trans.\ Robotics}, 33(1), 1--21.
\item Kaess, M., \emph{et al.} (2012). iSAM2: Incremental smoothing and mapping using the Bayes tree. \emph{Int.\ J.\ Robotics Research}, 31(2).
\item Besl, P.\ J., \& McKay, N.\ D. (1992). A method for registration of 3-D shapes. \emph{IEEE PAMI}, 14(2), 239--256.
\item Segal, A., Haehnel, D., \& Thrun, S. (2009). Generalized-ICP. \emph{Robotics: Science and Systems}.
\item Biber, P., \& Stra{\ss}er, W. (2003). The normal distributions transform: a new approach to laser scan matching. \emph{IROS}.
\item Arun, K.\ S., Huang, T.\ S., \& Blostein, S.\ D. (1987). Least-squares fitting of two 3-D point sets. \emph{IEEE PAMI}, 9(5).
\item Thrun, S., Burgard, W., \& Fox, D. (2005). \emph{Probabilistic Robotics}. MIT Press.
\item Brenner, C. SLAM lecture series (Units A--G), Leibniz University Hannover (video lectures and Python exercises).
\item ASPRS (2023; Version 2, 2024). \emph{ASPRS Positional Accuracy Standards for Digital Geospatial Data, Edition 2}. American Society for Photogrammetry and Remote Sensing.
\item Florida Administrative Code, Chapter 5J-17, \emph{Standards of Practice} (Board of Professional Surveyors and Mappers). Consult the current text.
\item Florida Department of Transportation. Terrestrial mobile LiDAR surveying and mapping guidance (current edition).
\item Groves, P.\ D. (2013). \emph{Principles of GNSS, Inertial, and Multisensor Integrated Navigation Systems}, 2nd ed. Artech House.
\item Dellaert, F., \& Kaess, M. (2017). Factor graphs for robot perception. \emph{Foundations and Trends in Robotics}, 6(1--2).
\item IEC 60825-1. Safety of laser products --- Part 1: Equipment classification and requirements.
\end{enumerate}

\vspace{1em}
\begin{bigidea}
\textbf{Where to go next.} Run LIO-SAM or a similar open-source pipeline on a public dataset; then take a commercial system to a small, well-controlled site, survey 30 checkpoints, and follow Module 21 from plan to accuracy statement. Nothing teaches like closing your own loop.
\end{bigidea}
